% TOP_PROBLEM: {"id": 30006874, "title": "Hopf Conjecture: no metric of positive sectional curvature on S^2 x S^2", "category_id": 6, "category": {"id": 6, "name": "geometry", "display_name": "Geometry", "description": "Euclidean and non-Euclidean geometry, geometric structures.", "slug": "geometry", "order_index": 6, "created_at": "2026-07-31T15:26:25.671Z"}, "status": "open"}
% FIRST_POSTED: 2026-09-27
% !TeX program = lualatex
% ATTEMPT_STATUS: unresolved
\documentclass[11pt]{article}
\usepackage[margin=25mm]{geometry}
\usepackage{fontspec}
\setmainfont{Latin Modern Roman}
\usepackage{amsmath,amssymb,mathtools}
\usepackage{unicode-math}
\setmathfont{Latin Modern Math}
\usepackage{xurl,hyperref}
\hypersetup{colorlinks=true,urlcolor=blue}
\setcounter{secnumdepth}{0}
\setlength{\parindent}{0pt}
\setlength{\parskip}{5pt}
\setlength{\emergencystretch}{3em}
\begin{document}
\section*{\texorpdfstring{Hopf Conjecture: no metric of positive sectional curvature on S\textasciicircum{}2 x S\textasciicircum{}2}{Hopf Conjecture: no metric of positive sectional curvature on S\textasciicircum{}2 x S\textasciicircum{}2}}\label{30006874}
\subsection{Definitions and mathematical statement}
Let $M=S^2\times S^2$ with its usual smooth structure. For a Riemannian metric $g$ and a two-dimensional tangent plane $\sigma=\operatorname{span}(u,v)$, sectional curvature is
\[
 \sec_g(\sigma)=\frac{\langle R(u,v)v,u\rangle_g}
 {|u|_g^2|v|_g^2-\langle u,v\rangle_g^2}.
\]
The Hopf assertion is that no smooth Riemannian metric on $M$ has $\sec_g(\sigma)>0$ at every point and every tangent two-plane. Thus the existence question permits arbitrary metrics, not only product metrics or metrics with symmetries. The standard product metric has flat mixed planes, but observing this for that particular metric does not rule out positive curvature for other metrics.
\par\smallskip\textit{Formulation and concept references:} \href{https://arxiv.org/html/math/0701389v3}{[S1]}.

\subsection{Short English statement}
Can the product of two two-spheres carry any Riemannian metric with positive sectional curvature in every tangent two-plane?
\subsection{Research attempt: restricted metric obstructions and a recent contrary claim}
\paragraph{1. The historical statement and the current source qualification.}
The preserved question is the unrestricted nonexistence assertion for positively curved metrics on the usual smooth $S^2\times S^2$. The survey [S1] records that historical conjecture. A current primary-source check on 27 September 2026 finds a materially different situation: Brendle and Hung's August 2026 preprint [E1], Theorem 1.1, asserts existence of such a metric. Its construction uses a third-order perturbation of a Cheeger--M\"uter metric, with supporting Mathematica calculations. This directly concerns the stated manifold and sectional curvature; it is not merely positive Ricci curvature or positivity on most planes. The theorem statement and the construction's scope were inspected, but its full analytic estimates and attached notebook were not independently certified in this attempt. Accordingly the historical assertion must not be advertised here as unconditionally still open, and the external counterexample claim must not be advertised as an independently verified construction supplied below.

Our own work gives rigorous obstructions for several special classes of metrics. Those obstructions are compatible with a metric outside the special classes. They do not refute [E1], nor do they prove the catalogue's universal negative assertion.

\paragraph{2. The exact reason the product metric fails.}
For a product metric $g=g_1\oplus g_2$ on $M_1\times M_2$, the Levi--Civita connection splits. If $X$ and $Y$ are vector fields lifted from different factors, then $\nabla_XY=\nabla_YX=0$, and the curvature tensor satisfies $R(X,Y)Y=0$. At any point, choose unit vectors $X\in TM_1$ and $Y\in TM_2$. Their span is a mixed two-plane, and
\[
 \sec_g(X\wedge Y)=\langle R(X,Y)Y,X\rangle=0.
\]
Thus no product metric on two positive-dimensional factors has strictly positive sectional curvature. This computation permits completely arbitrary metrics on the individual spheres. The topology of a product, however, does not force an arbitrary metric to have a parallel product splitting. There is no justified passage from ``every product metric fails'' to ``every metric on this product fails.''

\paragraph{3. Every global conformal deformation of a product still fails.}
Let $M_1,M_2$ be closed positive-dimensional manifolds, let $g=g_1\oplus g_2$, and let $\widehat g=e^{2u}g$ for a smooth real function $u$. For $g$-orthonormal $X,Y$, the conformal sectional-curvature formula is
\[
 \widehat K(X,Y)=e^{-2u}\bigl(K(X,Y)
 -\operatorname{Hess}u(X,X)-\operatorname{Hess}u(Y,Y)
 +du(X)^2+du(Y)^2-|\nabla u|^2\bigr).
\]
For completeness, this formula follows by substituting
\[
 \widehat\nabla_XY=\nabla_XY+du(X)Y+du(Y)X-g(X,Y)\nabla u
\]
into the definition of curvature and evaluating on an orthonormal pair. At a global minimum $p$ of $u$, one has $du(p)=0$ and $\operatorname{Hess}_p u\geq0$. Take a mixed plane at $p$. Its original sectional curvature is zero, so the formula yields
\[
 \widehat K_p(X,Y)
 =-e^{-2u(p)}\bigl(\operatorname{Hess}_p u(X,X)
 +\operatorname{Hess}_p u(Y,Y)\bigr)\leq0.
\]
Therefore no metric globally conformal to a product metric on a closed product can have positive sectional curvature everywhere. Compactness supplies the minimum. The statement is stronger than the mixed-plane calculation but remains a restriction on a particular family of metrics. In particular, a general symmetric two-tensor perturbation is not necessarily of the form $2u g$.

This calculation also checks the direction of the extremum. Using a maximum would reverse the useful Hessian sign and would not give the obstruction. The accompanying algebraic diagnostic checks the curvature expression with zero gradient and a positive semidefinite diagonal Hessian; its purpose is a sign check, not numerical evidence about all metrics.

\paragraph{4. Warping one factor cannot help either.}
Consider a warped product
\[
 g=g_B+f(x)^2g_F
\]
on a closed base $B$ and positive-dimensional fiber $F$, where $f>0$ is smooth. Let $X$ be a base unit vector and let $U$ be a fiber vector normalized to be unit for $g$. The connection identities, obtained from the Koszul formula, include
\[
 \nabla_XV=X(\log f)V,
 \qquad (\nabla_VW)^{\mathrm{horizontal}}
 =-g(V,W)\nabla\log f
\]
for lifted vertical fields. Substitution into the curvature tensor gives the mixed-plane identity
\[
 \sec_g(X\wedge U)=-\frac{\operatorname{Hess}_B f(X,X)}{f}.
\]
At a minimum of $f$, the right side is nonpositive for every $X$. Thus there is no positively curved warped product with closed base and a nontrivial fiber. Constant $f$ recovers the product case. This proof does not assume a round base or fiber, and it does not use the curvature of either factor.

The missing freedom is now explicit: a hypothetical unrestricted metric may have cross terms, nonintegrable horizontal distributions, and fiber geometry that varies in ways not encoded by a single warping function. Removing those possibilities without proof would simply insert the desired conclusion into the ansatz.

\paragraph{5. A global obstruction when two fibers remain totally geodesic.}
A useful form of Frankel's intersection argument can be proved directly. Let $P,Q$ be disjoint closed totally geodesic submanifolds of a closed positively curved $n$-manifold, with $\dim P+\dim Q\geq n$. Choose a shortest unit-speed geodesic $\gamma:[0,L]\to M$ from $P$ to $Q$. It meets both submanifolds orthogonally. Parallel transport $T_{\gamma(0)}P$ to the endpoint. Both it and $T_{\gamma(L)}Q$ lie in the $(n-1)$-dimensional orthogonal complement of $\dot\gamma(L)$, so their intersection has positive dimension. A nonzero vector in that intersection extends to a parallel normal field $V$ along $\gamma$ with endpoints tangent to $P$ and $Q$.

The second variation for curves whose endpoints lie on $P,Q$ has zero endpoint second-fundamental-form terms because both submanifolds are totally geodesic. Its index form is
\[
 I(V,V)=\int_0^L\left(|D_tV|^2
 -\langle R(V,\dot\gamma)\dot\gamma,V\rangle\right)dt<0.
\]
The derivative term vanishes and the curvature term is strictly positive. This contradicts the minimizing property, which requires nonnegative second variation for every admissible field. Hence such $P,Q$ must intersect.

For $S^2\times S^2$, two distinct fibers of either coordinate projection are disjoint closed surfaces and have dimensions summing to four. Therefore a positively curved metric cannot make both of these surfaces totally geodesic. This is a genuine obstruction if total geodesicity has been independently established. The smooth fibers exist for every metric, but their second fundamental forms generally do not vanish. In that case the omitted boundary terms have no controlled sign, and the argument stops.

\paragraph{6. Positive sectional curvature is weaker than a positive curvature operator.}
A common tempting route is to use $b_2(S^2\times S^2)=2$ and a Bochner vanishing theorem for a positive curvature operator. Its curvature hypothesis is stronger than the one in the question. The distinction can already be seen in an algebraic curvature tensor in dimension four.

Choose an orientation and decompose $\Lambda^2=\Lambda^+\oplus\Lambda^-$. Consider the symmetric operator with blocks
\[
 \mathcal R=\begin{pmatrix}A&0\\0&C\end{pmatrix},
 \qquad A=\operatorname{diag}(-1/2,7/4,7/4),\quad C=I_3.
\]
The traces of $A$ and $C$ are both three, which is precisely the first-Bianchi trace condition for such a block-diagonal symmetric operator in dimension four. Thus it represents an algebraic curvature tensor. A unit simple two-form has the form $(u+v)/\sqrt2$ with $u\in\Lambda^+$ and $v\in\Lambda^-$ unit. Its sectional curvature is
\[
 \langle\mathcal R\omega,\omega\rangle
 =\tfrac12(\langle Au,u\rangle+\langle Cv,v\rangle)
 \geq\tfrac12(-1/2+1)=1/4>0.
\]
Nevertheless $\mathcal R$ has a negative eigenvalue. This is a pointwise algebraic example, not a metric on the target manifold. It proves that replacing sectional positivity by operator positivity is not a valid pointwise implication. The ordinary product of round spheres already has positive Ricci and scalar curvature, so those weaker conditions cannot supply the desired obstruction either.

\paragraph{7. Perturbation diagnostics and the precise limit of this attempt.}
Even if a perturbation makes every previously flat plane positive at each fixed point for sufficiently small parameter, the choice of parameter must work uniformly over the compact Grassmann bundle. Pointwise thresholds can tend to zero near a degenerate stratum. A model warning is $f(x,s)=x^2s-s^2$ on $[0,1]$: each fixed $x>0$ permits positive values for sufficiently small $s>0$, while $f(0,s)<0$. A complete positive-curvature construction must handle such degeneracy; a complete obstruction must also allow the minimizing plane to move with the metric. Reading only the first variation at the old flat plane does neither.

The primary source [E1] explicitly addresses higher-order perturbations and degenerate loci. The elementary restrictions established here do not audit those estimates. No Mathematica execution or global curvature certificate for [E1] is claimed by this document.

\textbf{UNRESOLVED as an independent attempt.} Product, warped-product, conformal-product, and two-totally-geodesic-fiber approaches are ruled out by complete arguments. The unrestricted historical nonexistence assertion is not proved; the current source audit records a directly contrary 2026 theorem claim. The remaining task for a verified status change is independent validation of that construction, rather than an extrapolation of any restricted obstruction above.

\subsection{Sources}
\begin{itemize}
\item[S1]W. Ziller, Examples of Riemannian Manifolds with non-negative sectional curvature, arXiv:math/0701389v3 [math.DG], 2007.
\url{https://arxiv.org/html/math/0701389v3}.
\textit{Formulation source. Consulted 24 September 2026: Section 1: Hopf curvature conjectures and Bott ellipticity over every coefficient field..}
\item[E1]S. Brendle and P.-K. Hung, A metric on $S^2\times S^2$ with positive sectional curvature, arXiv:2608.19068v2, 24 August 2026, Theorem 1.1 and introduction.
\url{https://arxiv.org/html/2608.19068v2}.
\textit{Primary current-status source checked 27 September 2026. The existence assertion and its scope were inspected; the complete perturbation proof and attached Mathematica code were not independently certified here.}
\end{itemize}
\end{document}
